\begin{frame}[fragile]{A partition packs blocks iff it cuts the image's \emph{inner} axis --- so the rule is a diagonal}

\begin{columns}[T]
\begin{column}{0.50\textwidth}
\footnotesize
LDS is linear, so the tile is serialised in one of two orders, chosen per operand by the
\key{memory layout} of the global tensor:

\smallskip
\begin{tabular}{@{}ll@{}}
\toprule
unroll-major & image \texttt{[free][unroll]} \\
tile-major   & image \texttt{[unroll][free]} \\
\bottomrule
\end{tabular}

\smallskip
The hardware writes each region \key{densely}. Cut the \emph{outer} axis and the regions are runs
of whole rows, adjacent --- the unsplit image is reproduced exactly. Cut the \emph{inner} axis and
a dense write cannot interleave stripes, so what lands is a different image.

\medskip
Stated once, in \texttt{tdm\_split.split\_packs\_lds}:

\begin{lstlisting}[style=py]
return (axis == AXIS_MT) == (not unrollMajor)
\end{lstlisting}
\tsub{an XNOR of "the cut is on the free axis" with "the free axis is inner"}
\end{column}

\begin{column}{0.48\textwidth}
\centering
\begin{tabular}{@{}lcc@{}}
\toprule
cut axis & unroll-major & tile-major \\
         & \tsub{[free][unroll]} & \tsub{[unroll][free]} \\
\midrule
\texttt{AXIS\_MT} \tsub{free}   & \good{contiguous} & \bad{PACKED} \\
\texttt{AXIS\_DU} \tsub{unroll} & \bad{PACKED}      & \good{contiguous} \\
\bottomrule
\end{tabular}

\medskip
\footnotesize
\key{Both cells of the MT row are measured --- one in the negative.}

\smallskip
\begin{tabular}{@{}p{0.94\textwidth}@{}}
\toprule
\bad{MT, tile-major $\rightarrow$ PACKED} (2026-08-17, NT MT32x32x64): without the packed
treatment every NT split kernel read region~0's bytes for half its tiles. \\
\midrule
\good{MT, unroll-major $\rightarrow$ CONTIGUOUS} (2026-08-16): \emph{adding} a region term here
regressed 82 passing kernels, 110 failures $\rightarrow$ 192. \\
\bottomrule
\end{tabular}

\smallskip
Same axis, two layouts, opposite treatment. \bad{Generalising either measurement to the other
layout is the defect.}
\end{column}
\end{columns}
\end{frame}
